problem:  
Let \((X,d,\mathfrak{m})\) be a compact, essentially non-branching metric measure space satisfying the curvature-dimension condition \(\mathrm{CD}(K,N)\) for some \(K>0\) and finite \(N>1\) in the Lott–Sturm–Villani sense.  

For any probability measures \(\mu_0 = \rho_0 \mathfrak{m}\) and \(\mu_1 = \rho_1 \mathfrak{m}\), where \(\rho_0,\rho_1 \in L^\infty(X,\mathfrak{m})\) have bounded support and finite Cheeger energy, let \((\mu_t^{\varepsilon})_{t\in[0,1]}\) denote the **entropic interpolation** (the marginal flow of the Schrödinger bridge) corresponding to the entropy-regularized optimal transport problem with regularization parameter \(\varepsilon > 0\). Assume that for each sequence \(\varepsilon_j \to 0\):  

1. The interpolations \(\mu_t^{\varepsilon_j}\) converge to the displacement interpolation \((\mu_t)\) in the \(L^2\)-Wasserstein topology.  
2. The densities \(\mu_t^{\varepsilon_j} = \rho_t^{\varepsilon_j} \mathfrak{m}\) remain absolutely continuous for all \(t\in[0,1]\).  
3. The functions \(\rho_t^{\varepsilon_j}\) are uniformly Lipschitz with respect to the metric \(d\), with a Lipschitz constant independent of \(t,\varepsilon_j\), and of the initial data within bounded \(L^\infty\)-norm.  

**Question:**  
Does the existence of such uniform Lipschitz regularity of entropic interpolations for all bounded, compactly supported pairs \((\rho_0,\rho_1)\) imply that \((X,d,\mathfrak{m})\) must be isomorphic (in the measured Gromov sense) to a smooth weighted Riemannian manifold \((M,g,e^{-f}\mathrm{dvol}_g)\) whose Bakry–Émery Ricci curvature satisfies  
\[
\mathrm{Ric}_g + \nabla^2 f \ge K g, \quad \dim M \le N\ ?
\]  
Equivalently, can the **uniform regularity of Schrödinger bridges** among CD(K,N) metric measure spaces serve as a new analytic characterization of smooth Riemannian structure?  

Why is it a "good" problem:  
This problem formalizes an analytic–geometric rigidity question at the frontier of optimal transport and synthetic geometry. It asks whether strong analytic regularity of entropic interpolations—an object from stochastic optimal transport—forces the underlying synthetic space to be smooth and Riemannian.  

1. **Cross-disciplinary unification:**  
   The problem connects entropic optimal transport (a probabilistic and analytic construct) with Lott–Sturm–Villani curvature bounds (a geometric and synthetic framework). Establishing such a characterization would conceptually link stochastic regularization with geometric smoothness.  

2. **Beyond known equivalences:**  
   In smooth settings, CD(K,N) and Bakry–Émery curvature lower bounds coincide. This problem reverses direction: it seeks to *recover smoothness* of the space via an analytic regularity property. No existing theorem provides such a criterion, making this direction open and deep.  

3. **Simplicity with depth:**  
   The statement is concise—does uniform Lipschitz regularity of entropic interpolations characterize Riemannian structure?—but resolving it would require combining PDE analysis on Wasserstein space, metric measure geometry, and curvature-dimension theory.  

4. **Potential for structural breakthroughs:**  
   A positive answer would yield an analytic test for identifying when a CD(K,N) space is Riemannian by purely transport-theoretic means, paralleling how the linearity of the heat flow characterizes RCD spaces. It would thus unify entropy-regularized transport theory with the geometric analysis foundations of modern curvature-dimension geometry.  

Hence, the problem sits at the interface of geometry, analysis, and probability, embodying the conceptual insight and potential impact expected of a truly “good” mathematical problem.